\begin{longtable}{@{}p{5cm}p{10cm}@{}}
\toprule
\textbf{Término} & \textbf{Definición} \\ \midrule
\endhead
Operación específica & Actividad con \gls{uas} que, por su nivel de riesgo o complejidad, requiere evaluación y mitigación operacional conforme al \gls{rca} 100. \\
\gls{uas} (Unmanned Aircraft System) & Sistema compuesto por la aeronave no tripulada, la estación de control, el enlace de comunicaciones y todos los elementos asociados para su operación segura. \\
Piloto & Persona calificada y certificada para operar un \gls{uas} desde una estación de control remota. \\
Observador & Miembro de la tripulación responsable de mantener contacto visual continuo con el \gls{uas} para apoyar la operación segura. \\
Vuelo \gls{vlos} & Operación en la que el piloto o un observador designado mantiene contacto visual directo con la aeronave. \\
Vuelo \gls{evlos}  & Operación en la que se mantiene el contacto visual con el UA, como en los \gls{vlos}, pero al doble de distancia porque se cuenta con la figura de un observador, que siempre se mantendrá en contacto por radio con el piloto. \\
Vuelo \gls {bvlos} & Operación en la que la aeronave vuela más allá del alcance visual directo, cumpliendo las mitigaciones y autorizaciones necesarias. \\
Vuelo \gls{fpv} & Operación en la que el piloto controla el \gls{uas} utilizando una vista en primera persona transmitida por una cámara a gafas o monitor, complementada con observadores visuales cuando sea necesario. \\
Zona de operación & Área geográfica definida para la ejecución de la misión, incluyendo perímetro de seguridad y puntos críticos. \\
\gls{metar} & Informe meteorológico aeronáutico que describe las condiciones del tiempo observadas en un aeródromo o zona de interés. \\
\gls{notam} & Aviso distribuido por medios de telecomunicaciones que contiene información sobre el establecimiento, condición o cambio en una instalación aeronáutica, servicio, procedimiento o peligro. \\
Mitigación operacional & Acción o procedimiento destinado a reducir el nivel de riesgo asociado a un peligro identificado. \\
Contingencia & Procedimiento planificado para responder a un evento inesperado que afecte la operación segura del \gls{uas}. \\
Operación sobre multitudes & Actividad en la que el \gls{uas} sobrevuela concentraciones de personas, siguiendo requisitos adicionales de seguridad y mitigación. \\
Transmisión en vivo (Broadcast) & Operación audiovisual en la que las imágenes captadas por el \gls{uas} se transmiten en tiempo real a una estación de producción o medio de comunicación. \\
Pérdida de enlace C2& Interrupción del enlace de control o video entre el \gls{uas} y la estación de control. \\
\bottomrule
\end{longtable}
